Can a single set of weights act as a ``computer''?  
We term this abstraction a Neural Computer (NC): a neural system that unifies computation, memory, and I/O in a learned runtime state.
This usage is distinct from the Neural Turing Machine / Differentiable Neural Computer line~\citep{graves2014neural,graves2016hybrid}: our concern is not differentiable external memory, but whether a learning machine can begin to assume the role of the running computer itself.

To implement this idea, we instantiate NCs as video models.
At this stage, video models are the most practical substrate for this prototype, though we expect the long-term solution to require a fundamentally new neural architecture (\Cref{section:toward-cnc}).
This implementation draws on several technical lines. World models~\citep{ha2018world} show that neural networks can internalize environment dynamics and support predictive imagination, while high-capacity video generators such as Veo~3.1~\citep{google_veo3_1_2025} and Sora~2~\citep{openai_sora2_2025} show that such learned dynamics can be rendered into coherent frame sequences.
Frontier interactive video models such as Genie~3~\citep{bruce2024genie} further extend this trajectory toward action-controllable generative environments. These lines provide practical machinery for current NC prototypes, but do not by themselves define the NC abstraction.
In parallel, LLM-driven UI systems such as Imagine with Claude\footnote{\url{https://claude.ai/imagine/}} map natural-language inputs to structured interface updates. 
Yet these capabilities remain split across different systems objects: conventional computers execute explicit programs, agents act through external execution environments, and world models render or predict environment dynamics, while executable state still resides outside the model.
NCs are motivated by this gap: they are not a smarter layer on top of the existing stack, but a proposal to make the model itself the running computer.
The immediate question in this paper is whether early runtime primitives can be learned directly from raw interface I/O without privileged access to program state.

Throughout this paper, NC denotes this proposed machine form, while CNC denotes its mature, general-purpose realization.
We study two interface-specific prototypes of this NC formulation (see~\Cref{section:prelim}).
\nccligen{} models terminal interaction from text (natural language or command lines) and an initial frame, while \ncguiworld{} models desktop interaction from recent pixels and synchronized mouse/keyboard actions (\Cref{section:impl-cligen,section:impl-guiworld}).

\begin{tcolorbox}[
  colback=apricotglaze!40!white,
  colframe=metablue!40,
  interior style={shade,shading angle=315,left color=white,right color=apricotglaze!55!white},
  boxrule=0pt,
  borderline west={1pt}{0pt}{gray!60},
  left=6pt,right=4pt,top=3pt,bottom=3pt]
\small
\textbf{Neural Computer (NC) abstraction~(\hyperlink{teaser}{Teaser}).}
A neural system $(F, G)$ parameterized by $\theta$ that models an interactive computer interface through a single latent runtime state $h_t$ that carries executable interface state and also acts as working memory (see~\cref{eq:wm}).
\end{tcolorbox}

In the \nccligen{} experiments, the NC can render and execute basic command-line workflows.
It often stays aligned with the terminal buffer and captures common ``physics'' of everyday CLI use (e.g., fast scrollback, prompt wrapping, window resizing).
On carefully scripted data, rollouts can be visually and structurally close to real sessions, and the model can execute short command chains and render their outputs.
Arithmetic-probe scores improve substantially with stronger system-level conditioning, though symbolic stability remains limited.

In the \ncguiworld{} experiments, we evaluate standard world-model designs across action injection, action encoding, and data quality.
Figure~\ref{fig:nc-general} summarizes this template across two interface-specific NCs trained separately without shared parameters.
Qualitatively, the model learns coherent pointer dynamics and short-horizon action responses (e.g., hover/click feedback and window/menu transitions), suggesting that local GUI control primitives are learnable in controlled settings.

Our experimental insights indicate that current NCs already realize early runtime primitives, most notably I/O alignment and short-horizon control.
The long-term target is a Completely Neural Computer (CNC), the mature, general-purpose realization of this machine form: a fully learned computer whose compute, memory, and interfaces are unified in a single learned runtime substrate rather than engineered as separate modules.
These prototypes are an early step toward that CNC vision. 
Substantial challenges remain in robust long-horizon reasoning, reliable symbolic processing, stable capability reuse, and explicit runtime governance.
\Cref{section:toward-cnc} outlines these open challenges and a roadmap toward CNCs. 

\begin{tcolorbox}[
  colback=apricotglaze!40!white,
  colframe=metablue!40,
  interior style={shade,shading angle=315,left color=white,right color=apricotglaze!55!white},
  boxrule=0pt,
  borderline west={1pt}{0pt}{gray!60},
  left=6pt,right=4pt,top=3pt,bottom=3pt]
\small
\textbf{Completely Neural Computer (CNC) abstraction~(\Cref{sec:roadmap}).}
A Neural Computer instance is \emph{complete} (i.e., a CNC) if it is
(i) \emph{Turing complete},
(ii) \emph{universally programmable},
(iii) \emph{behavior-consistent} unless explicitly reprogrammed, and
(iv) realizes the architectural and programming-language advantages of NCs relative to conventional computers.
\end{tcolorbox}

\textbf{Concretely, this work makes the following contributions:}
{
\renewcommand\labelitemi{\large$\bullet$}
\begin{itemize}
\setlength{\itemsep}{2pt plus 1pt minus 1pt}
\setlength{\parskip}{0pt}
\item Define neural computers (NCs) and build video-based prototypes for both CLI and GUI interfaces.
\item Provide a data engine and alignment recipe that synchronize text, actions, and frames for the CLI and GUI environments used in this paper.
\item Identify practical design choices for NCs through extensive ablation studies.
\item Outline an engineering roadmap toward completely neural computers (CNCs), centered on acceptance challenges such as reuse, consistency, and runtime governance.

\end{itemize}
}
